\clearpage
\VerdierPhysicalPageAnchor{0132}%
\begingroup
\small
\setlength{\emergencystretch}{1.5em}%
\setlength{\abovedisplayskip}{2pt plus 1pt minus 1pt}%
\setlength{\belowdisplayskip}{2pt plus 1pt minus 1pt}%

\VerdierSemanticAnchor{ch2-2-1-10-continuation}
\noindent SM3) が従う。
\VerdierSemanticAnchor{ch2-2-1-10-triangle-setup}%
仮定により、$N$ が $\mathcal D'$ の対象であるような完全三角
$(X,Y,N,s,h,u)$ が存在する。$fs=0$ なので、(1.2.1) により
$f=gh$ を満たす射 $g\colon N\to Z$ が存在する。

\VerdierSemanticAnchor{ch2-2-1-10-factorization}
次に $g$ を含む完全三角 $(N,Z,W,g,t,v)$ をとる (TRI)。このとき
$tg=0$ (1.2.2)、したがって $tf=tgh=0$ である。

\VerdierSemanticAnchor{ch2-2-1-10-closure-conclusion}
よって $t$ が $S(\mathcal D')$ に属することを示せばよく、これは TRII)
から直ちに従う。

\VerdierSemanticAnchor{ch2-2-1-11}
\medskip
\noindent\textbf{2.1.11.}\quad
\VerdierSemanticAnchor{ch2-2-1-11-sm5}
\emph{$S(\mathcal D')$ は性質} SM5) \emph{をもつ。}
\VerdierSemanticAnchor{ch2-2-1-11-proof}
これは TRII) と $\mathcal D'$ がシフトで安定であることから直ちに従う。

\VerdierSemanticAnchor{ch2-2-1-12}
\medskip
\noindent\textbf{2.1.12.}\quad
\VerdierSemanticAnchor{ch2-2-1-12-sm6}
\emph{$S(\mathcal D')$ は性質} SM6) \emph{をもつ。}
\VerdierSemanticAnchor{ch2-2-1-12-triangle-pair}
$(X,Y,Z,u,v,w)$ と $(X',Y',Z',u',v',w')$ を $\mathcal D$ の二つの
完全三角とし、

\VerdierSemanticAnchor{ch2-2-1-12-input-square}
\begin{center}
  \begin{tikzcd}[column sep=4.2em,row sep=2.5em]
    X \arrow[r,"u"] \arrow[d,"s"'] &
    Y \arrow[d,"{s'}"] \\
    X' \arrow[r,"{u'}"] & Y'
  \end{tikzcd}
\end{center}

\VerdierSemanticAnchor{ch2-2-1-12-completion-goal}
\noindent を、$s$ と $s'$ が $S(\mathcal D')$ に属する可換図式とする。
$(s,s',s'')$ が三角の射となるような $S(\mathcal D')$ の射
$s''\colon Z\to Z'$ を見つけなければならない。

\VerdierSemanticAnchor{ch2-2-1-12-case-reduction}
まず次の二つの場合のいずれかにあると仮定してよい：

\begin{enumerate}
\renewcommand{\labelenumi}{\emph{\alph{enumi})}}
  \item \VerdierSemanticAnchor{ch2-2-1-12-case-a}%
  $X=X'$ かつ $s=\operatorname{id}_X$。
  \item \VerdierSemanticAnchor{ch2-2-1-12-case-b}%
  $Y=Y'$ かつ $s'=\operatorname{id}_Y$。
\end{enumerate}

\VerdierSemanticAnchor{ch2-2-1-12-factorization}
実際、任意の射の射は次のように分解する：

\VerdierSemanticAnchor{ch2-2-1-12-factorization-diagram}
\begin{center}
  \begin{tikzcd}[column sep=4.2em,row sep=1.7em]
    X \arrow[r,"u"] \arrow[d,"\operatorname{id}_X"'] &
    Y \arrow[d,"{s'}"] \\
    X \arrow[r,"{s'u}"] \arrow[d,"s"'] &
    Y' \arrow[d,"\operatorname{id}_{Y'}"] \\
    X' \arrow[r,"{u'}"] & Y'
  \end{tikzcd}
  \quad ,
\end{center}

\VerdierSemanticAnchor{ch2-2-1-12-composition}
\noindent そして $s''$ は合成により得られる（$S(\mathcal D')$ が合成について
閉じていることはすでに示した）。

\VerdierSemanticAnchor{ch2-2-1-12-opposite-category}
次に、反対圏へ移ることで場合 b) は場合 a) から従う。

\VerdierSemanticAnchor{ch2-2-1-12-triv-conclusion}
したがって場合 a) だけを扱えばよく、これは TRIV) から直ちに得られる帰結
である。

\VerdierSemanticAnchor{ch2-2-1-13-p133-boundary}
\endgroup
